\chapter{Warning Signs During Pregnancy}

\section*{Definition}
A set of signs and symptoms during pregnancy that may indicate serious complications requiring urgent medical evaluation, including preeclampsia, placental abruption, preterm labor, or infection.

\section*{Pathophysiology}
Pregnancy induces significant physiological changes in the cardiovascular, renal, respiratory, and immune systems. Warning signs reflect pathological decompensation: hypertension and proteinuria indicate preeclampsia (endothelial dysfunction, placental ischemia); vaginal bleeding suggests placental abruption or placenta previa; decreased fetal movement indicates fetal distress; severe abdominal pain may signal placental abruption, uterine rupture, or appendicitis; fever suggests chorioamnionitis or other infection.

\section*{Etiology}
\begin{itemize}
  \item Preeclampsia (hypertension, proteinuria, headache, visual changes, RUQ pain)
  \item Placental abruption (vaginal bleeding, abdominal pain, uterine tenderness)
  \item Placenta previa (painless vaginal bleeding in third trimester)
  \item Preterm labor (regular contractions, pelvic pressure, back pain before 37 weeks)
  \item Chorioamnionitis (fever, uterine tenderness, foul amniotic fluid)
  \item Gestational diabetes complications (polyuria, polydipsia, blurred vision)
  \item Deep vein thrombosis (unilateral leg swelling, pain, warmth)
  \item Pulmonary embolism (sudden shortness of breath, chest pain, hemoptysis)
  \item Miscarriage (vaginal bleeding, cramping in first trimester)
  \item Ectopic pregnancy (sharp abdominal/pelvic pain, shoulder pain, fainting)
\end{itemize}

\section*{Indications for Evaluation}
Seek care if:
\begin{itemize}
  \item Severe or worsening symptoms
  \item Any vaginal bleeding in pregnancy, severe headache with visual changes, sudden swelling of face/hands, decreased fetal movement, rupture of membranes, or signs of preterm labor before 37 weeks
  \item Accompanied by other concerning signs
\end{itemize}

\section*{Related Symptoms}
\begin{itemize}
  \item Abdominal pain
  \item Vaginal bleeding
  \item Headache
  \item Swelling
  \item Fever
\end{itemize}
\textbf{Note:} Always consider serious underlying conditions when evaluating persistent warning signs during pregnancy.

\section*{Self-Care / Home Management}
\begin{itemize}
  \item Rest and avoid activities that may aggravate the condition.
  \item Maintain adequate hydration and a balanced diet.
  \item Over-the-counter remedies may provide symptomatic relief (consult a pharmacist or doctor).
  \item Monitor symptoms and seek professional advice if they persist or worsen.
  \item Avoid self-medicating with prescription medications without medical guidance.
\end{itemize}

\section*{Diagnostic Approach}
\begin{itemize}
  \item A thorough history and physical examination are the first steps in evaluation.
  \item Laboratory tests (blood work, urinalysis) may be ordered based on clinical suspicion.
  \item Imaging studies (X-ray, ultrasound, CT, MRI) may be indicated when structural pathology is suspected.
  \item Specialist referral is arranged when the diagnosis remains uncertain or requires specific expertise.
\end{itemize}

\section*{Treatment / Management Overview}
\begin{itemize}
  \item Treatment targets the underlying cause identified through diagnostic evaluation.
  \item Supportive care and symptom management are provided while awaiting definitive therapy.
  \item Medications, lifestyle modifications, and physical therapies may be prescribed as appropriate.
  \item Follow-up ensures treatment effectiveness and allows timely adjustments to the care plan.
\end{itemize}

\clearpage